\headline{{\textbf{\Large\sffamily Seasonal Care Tips:}}\\
          {\textbf{\large\sffamily Mid\--May to Mid\--June}}}
\label{2025-05-care}
\noindent
From mid\--May to mid\--June, bonsai trees in South England enter a period of exuberant growth.
With daylight hours at their peak and temperatures consistently warm, trees are in full swing---developing foliage, extending shoots, and storing energy.
However, this time of abundance also demands careful attention.
Growth must be refined, pests controlled, and preparations made for the heat of summer.
Here's how to support your trees during this crucial stretch of the growing season.

\medskip
\topic{Managing and Refining Growth}
\noindent
By this stage, most deciduous bonsai---such as maples, hornbeams, and elms---will have produced strong spring growth.
Continued shoot pinching and light pruning help maintain the tree's silhouette and encourage finer ramification.
Where growth is overly vigorous, partial defoliation can be employed to improve light penetration and balance energy distribution within the tree.

Pines may now be ready for decandling, particularly black and red pines.
This technique, best done by mid\--June, encourages back budding and a second flush of shorter, more compact needles.
Junipers benefit from careful trimming, focusing on maintaining foliage pads while avoiding excessive thinning, which can stress the tree.

\medskip
\topic{Watering and Nutrient Management}
\noindent
With longer days and increased evaporation, watering becomes a key part of your daily routine.
Check soil moisture each morning, and water only when needed---ensuring a deep, thorough soak rather than a quick surface wetting.
Avoid waterlogging by ensuring pots have good drainage and avoid trays that collect stagnant water.

Feeding should remain consistent to support sustained growth.
Use a balanced fertiliser weekly or apply slow\--release pellets suitable for your species.
For deciduous trees, a slightly nitrogen\--rich feed helps fuel foliage development, while for conifers and flowering bonsai, a more balanced or phosphorus\--heavy mix can enhance root and flower health.
Organic options such as seaweed extract or fish emulsion can be used alternately to promote overall vitality.

\medskip
\topic{Wiring, Shaping, and Maintenance}
\noindent
Branch flexibility continues to be favourable, making this a good time to wire and shape your trees.
Take care not to wire over tender new growth, particularly on evergreens, as it can bruise or break easily.
Review any existing wiring, as swelling branches may begin to bite into the bark---remove or adjust as necessary.

In addition to wiring, consider thinning overcrowded areas to improve airflow and light distribution.
Clean out the interior of dense canopies by removing small, shaded shoots that are unlikely to thrive.
Keep dead leaves and fallen debris out of pots to discourage pests and fungal issues.
\columnbreak

\medskip
\topic{Pests, Diseases, and Preventative Care}
\noindent
This is prime time for pests such as aphids, spider mites, vine weevils, and scale.
Look for sticky residue, distorted leaves, and fine webbing---common signs of infestation.
Treat early with insecticidal soap or neem oil, repeating weekly if needed.
Manual removal of pests with a soft brush or jet of water can also be effective.

Fungal diseases such as black spot, powdery mildew, and rust may emerge during damp or humid spells.
Ensure your trees have adequate space and airflow, and avoid watering foliage late in the day.
If you notice symptoms, remove infected leaves and consider using a mild fungicide.
Preventative care goes a long way---regular observation is your best defence.

\medskip
\topic{Sunlight, Shelter, and Positioning}
\noindent
With temperatures now regularly warm and the risk of frost long gone, most bonsai can be kept in full sun.
This promotes strong growth and helps harden new shoots.
However, on particularly hot days or during heatwaves, young trees or recently repotted specimens may benefit from partial shade during the hottest hours---especially species with tender foliage like Japanese maples or azaleas.

Tropical bonsai such as ficus, bougainvillea, and Chinese elm enjoy being outside full\--time now.
Place them in a sheltered, sunny spot protected from wind and sudden temperature drops.
Be mindful of night temperatures dipping below 10°C for more sensitive species.

\medskip
\topic{Looking Ahead: Summer Preparation}
\noindent
Mid\--May to mid\--June is the perfect time to begin preparing for summer.
As temperatures rise, consider setting up shade cloths or repositioning your trees to avoid excessive midday exposure.
Check irrigation systems or timers if you'll be away---test them thoroughly to prevent mishaps.

Moss or mulch can be added to help retain moisture and insulate roots, particularly in shallow pots.
Clean benches, tools, and the surrounding area to discourage pests and maintain good hygiene.
Now is also a good time to check pot stability and wiring anchors before winds or storms become more frequent.

\medskip
\topic{Taking Pleasure in the Season}
\noindent
This is one of the most rewarding periods in the bonsai calendar.
Trees are vibrant, full of life, and responding actively to your care.
Take a moment to enjoy the soft green of fresh needles, the architectural lines of newly refined branches, and the overall energy of your collection.
Photograph your trees to track progress and share updates with fellow enthusiasts---there's much to celebrate and learn together as the season unfolds.
\closearticle
%\columnbreak
